% -----------------------------------------------------------------------------
% main_v2.tex -- second draft of the Tethys data paper.
%
% This file is a ground-up rewrite of main.tex. The v1 draft is preserved at
% main.tex so coauthors can diff. Companion documents that inform this draft:
%   - EQUATIONS_AUDIT.md  -- per-equation reconciliation of paper vs code
%   - STYLE_NOTES.md      -- distilled voice/structure from 3 Bracken SciData papers
%
% Open issues flagged for coauthor resolution:
%   - Author list ordering and affiliation numbering (v1 had [3]/[4] mixed; I
%     kept the same names and reset the numbering here -- please verify).
%   - TGW-WRF primary reference is missing from Tethys.bib; marked with TODO.
%   - Bracken et al. 2025 GODEEEP-hydro reference also missing; marked with TODO.
% ------------------------------------------ÅÎÅ-----------------------------------

\documentclass[fleqn,10pt]{wlscirep}
\usepackage[utf8]{inputenc}
\usepackage[T1]{fontenc}
\usepackage{lineno}
\usepackage{booktabs}
\usepackage{easyReview}

\linenumbers

% \title{High-resolution 21st-century monthly water demands for the contiguous United States}
\title{High-resolution monthly sectoral water demands for the U.S. over 1980-2100} % across diverse futures

\author[1*]{Cameron Bracken}
\author[2*]{Hassan Niazi}
\author[1]{Travis Thurber}
\author[2]{Isaac Thompson}
\author[1]{Kazi Tamaddun}
\author[3]{Hisham Eldardiry}
\author[1]{Kendall Mongird}
\author[1,3]{Nathalie Voisin}
\author[1]{Ning Sun}
\author[1]{Jennie Rice}

\affil[1]{Pacific Northwest National Laboratory, Richland, WA, USA}
\affil[2]{Joint Global Change Research Institute, Pacific Northwest National Laboratory, College Park, MD, USA}
\affil[3]{University of Washington, Seattle, WA, USA}

\affil[*]{corresponding author(s): Cameron Bracken (cameron.bracken@pnnl.gov), Hassan Niazi (hassan.niazi@pnnl.gov)}

\begin{abstract}
% COMMENT Too much detail with little context for an opening. Needs an "agenda-setting" hook
Water demand in the United States totals roughly 322~Bgal~day$^{-1}$ across irrigation, thermoelectric, public-supply, industrial, and livestock sectors, with strong regional heterogeneity driven by climate, land use, population, and economic activity. Water demand projections that characterize the seasonal, spatial, and sectoral evolution of demands at sufficient resolution to capture local heterogeneity are essential for assessing sectoral water scarcity, quantifying cross-sectoral tradeoffs, and supporting integrated energy-water-land modeling. Yet most high-resolution demand datasets are single-sector, historical only. Future water demand projections typically inherit the coarse spatial and temporal scales of the integrated assessment models that produce them, and rarely capture the coupled evolution of energy, water, and land systems that jointly drive future demands. Here, we present a gridded, monthly, multi-sector dataset of water withdrawals and consumption for the contiguous United States (CONUS) at 1/8$^{\circ}$ (0.125$^{\circ}$) spatial resolution, covering 1980--2019 (historical) and 2020--2099 (eight future scenarios). The dataset covers irrigation, thermoelectric, municipal (public-supply and domestic), livestock, manufacturing, and mining demands, separately for withdrawals and consumption, and includes per-cell renewable vs. non-renewable source attributions. 
% COMMENT This is the language used in recent IM3/GCIMS papers to describe these scenarios. Would be good to keep the wording consistent.
Future scenarios represent a wide yet plausible range of atmospheric conditions, emissions constraints, and economic, technological, and population growth assumptions. We validate the historical record against 2015--2020 USGS water-use estimates at the HUC6 scale, finding agreement in annual total demand within 10\% at CONUS scale and correlations between 0.71 and 0.95 at HUC6. This dataset extends and improves the previously published Tethys global product\cite{Khan2023} by incorporating state-resolved, sectoral water demands and source attribution from the Global Change Analysis Model with subnational detail for the U.S. (GCAM-USA), explicit future power plant siting, scenario-consistent high-resolution atmosphere, land, and population projections, and USGS-based validation. 
\end{abstract}
\begin{document}

\flushbottom
\maketitle
% Click the title above to edit the author information and abstract

\thispagestyle{empty}

% -----------------------------------------------------------------------------
\section*{Background \& Summary}


\comment{Humans depend on water for irrigation, thermoelectric cooling, public supply, industry, and livestock, and the relative importance of these sectors varies sharply across regions and over time. Global water demand has continued to rise through the 21st century under combined socioeconomic and climatic pressures\cite{WWDR2019, Niazi2024PeakWater, Awais2024Preprint}, and demand-side drivers, not supply-side limits, dominate most projected shifts in water scarcity\cite{Graham_2020, Awais2024Preprint, Kyle2023Sustainability}. }{overall, i feel the review of both existing datasets and water scarcity modeling needs to be expanded significantly to more strongly motivate the need for this dataset. }
% COMMENT Good point, but this is a modeling issue, right? and would be better moved to support why higher-res demand matter for improved local scarcity modeling. And, this paragraph moves too quickly from water demand general statements to the modeling issue, missing a smooth transition.
Even where regional supply appears ample, scarcity can manifest locally when demand and availability are mismatched in space or time\cite{HADJIMICHAEL23}. Resolving water demand at the spatial and temporal scales of management decisions is therefore essential for evaluating scarcity under climate and socioeconomic change.

% COMMENT before presenting modeling challenges due to the lack of such data, suggest adding a paragraph reviewing existing water demand datasets in Table 1 and clarify their common limitations and why addressing these limitations matters in the context of water modeling, regarding spatial resolution, sectoral representation, and scenario plausibility. 
The models that can simulate water demand over long horizons and wide extents --- integrated human-Earth systems models such as the Global Change Analysis Model (GCAM)\cite{Zhao_gcamusa_water, Calvin2019GCAM, Binsted_2022} --- operate at coarse spatial units (typically 32 geopolitical regions or 235 basins globally) and 5-year time steps. Conversely, hydrology and water-management models that can represent fine-scale physical processes typically use simplified or static assumptions about the human water demand they route and supply\cite{Abeshu_2023, Wada_2017}. Bridging the two requires spatial and temporal downscaling of integrated-assessment outputs\cite{hess-17-4555-2013, van_Vliet_2021, Jones_2024}, a process in which the choice of gridded proxy variables and scenario-consistent forcing materially affects the resulting demand fields. Khan et al.\cite{Khan2023} first produced a global Tethys-downscaled multi-sector product at 1/2$^{\circ}$ resolution; subsequent work has improved scenario consistency in irrigation by coupling land-use downscaling to Demeter\cite{Vernon-2018, BIJL201675} and population downscaling to SSP projections\cite{Jones_2016}. 
% this feels too granular for the background section.. more important to clarify why these improvements are important or what gaps are filled through this list of improvements
No publicly available dataset, however, has yet combined GCAM-USA state-level demands with explicit power-plant siting, SSP-aware population, and climate-consistent forcing at a sub-state grid with monthly resolution across CONUS.

% COMMENT Suggest adding why monthly resolution is important, e.g., for capturing seasonal patterns in water demand that differ by sector and so on.
We present such a dataset here. The published record contains gridded monthly water withdrawals and consumption for six sectors --- irrigation, thermoelectric, municipal (public-supply and domestic), livestock, manufacturing, and mining --- at 1/8$^{\circ}$ spatial resolution, on the CONUS extent $[25.0625^{\circ}\text{N},\,52.9375^{\circ}\text{N}] \times [-124.9375^{\circ}\text{W},\,-67.0625^{\circ}\text{W}]$, for one historical period (1980--2019) and eight future scenarios (2020--2100) that vary by RCP, TGW climate sample, and SSP. For each scenario the dataset also provides per-cell renewable vs. non-renewable source shares. The downscaling chain improves on the prior Tethys CONUS product in six specific ways (Section~``Improvements over previous version''): GCAM-USA integration, explicit CERF-based\cite{Vernon2021} power-plant siting, SSP-aware population proxies, climate-forced irrigation temporal downscaling using TGW-WRF derived deficits and growing-season indices, USGS-anchored source-share adjustment, and resolution refinement from 1/2$^{\circ}$ to 1/8$^{\circ}$. Figure~\ref{fig:schematic} shows the overall workflow. Table~\ref{tab:prior-datasets} compares this dataset with closely related published records.

% COMMENT Suggest adding a box labeled Tethys to encompass the two green boxes. For temporal downscaling, change it to 5-year to monthly. In the Demeter box, add the downscaling scale: regional, 5-year to 1/8° spatial resolution, annual temporal resolution. Also spell out PET, HDD, and CDD in the figure caption, and "Monthly Climate Variables" should be "Monthly Meteorological Variables".
\begin{figure}[ht]
\centering
% \includegraphics[width=0.85\linewidth]{flow-chart}
\includegraphics[width=0.85\linewidth]{flow-chart2.pdf}
\caption{Workflow for producing the Tethys CONUS multi-sector water-demand dataset. HUC2 Basin-scale demands from GCAM-USA are spatially downscaled using sector-specific gridded proxies (Demeter land use, CERF power-plant locations, SSP population, and fixed livestock distributions), temporally downscaled using monthly meteorological indicators derived from the TGW meteorological dataset, and post-processed with USGS-anchored adjustment of per-cell renewable vs. non-renewable source shares.}
\label{fig:schematic}
\end{figure}

\begin{table}[ht]
\centering
\small
\begin{tabular}{lllll}
\toprule
\textbf{Dataset} & \textbf{Spatial res.} & \textbf{Temporal res.} & \textbf{Scenarios} & \textbf{Sectors} \\
\midrule
Huang et al.\cite{hess-22-2117-2018}           & 0.5$^{\circ}$ global & monthly & historical (1971--2010) & 6 sectors \\
Khan et al.\cite{Khan2023} (prior Tethys)      & 0.5$^{\circ}$ global & monthly & historical + SSP/RCP futures & 6 sectors \\
van Vliet et al.\cite{van_Vliet_2021}          & country/basin & annual & historical + SSP & 6 sectors \\
Wada et al.\cite{Wada_2017}                    & 0.5$^{\circ}$ global & monthly & historical + projections & 4 sectors \\
\textbf{This work}                             & \textbf{0.125$^{\circ}$ CONUS} & \textbf{monthly} & \textbf{hist.\ + 8 RCP/climate/SSP} & \textbf{6 sectors + GW/SW split} \\
\bottomrule
\end{tabular}
\caption{Gridded multi-sector water-demand datasets that overlap in scope with the product presented here. The current dataset increases spatial resolution fourfold over its immediate Tethys predecessor\cite{Khan2023}, adds per-cell groundwater/surface-water attribution, and uses scenario-consistent population, land-use, and climate forcing rather than static baselines.}
\label{tab:prior-datasets}
\end{table}

% -----------------------------------------------------------------------------
\section*{Methods and Data}

\subsection*{GCAM-USA}

Region-scale water-demand inputs come from the Global Change Analysis Model (GCAM-USA version)\cite{Zhao_gcamusa_water, Zhao2024, Calvin2019GCAM, Binsted_2022}. GCAM is a market-equilibrium integrated assessment model that allocates supply and demand across coupled energy, water, land, and economic sectors given scenario-specific assumptions on population, productivity, technology, and policy. GCAM-USA resolves U.S.~electricity generation, manufacturing, and municipal demands at the state level and crop-level irrigation demands at state $\times$ basin intersections, at 5-year intervals. 
% COMMENT Refer to the CERF paper I shared for scenario development and descriptions.
The eight GCAM-USA runs used in this study combine two RCPs (4.5, 8.5), two climate samples (cooler, hotter) drawn from the TGW perturbed-thermodynamics ensemble, and two SSPs (3, 5).

\subsection*{Preprocessing of meteorological forcing}

Irrigation and electricity temporal downscaling both require monthly meteorological variables consistent with each scenario. We derive these from the Thermodynamic Global Warming (TGW) meteorological forcing \cite{Jones2023TGW}at 1/8$^{\circ}$ covering the CONUS for the historical period and four future scnearios. Historical forcings were dynamically downscaled from ~30-km ERA5 reanalysis (1980-2019) using WRF. Future meteorological forcings were projected through the TGW approach, which replays historical weather sequences with added thermodynamic signals derived from CMIP6 Earth System Models. The TGW method simulates how historical weather events may shift under warming conditions while retaining the physical realism and synoptic structure of observed weather. Four scenarios combine diverging emission pathways and structural uncertainties in CMIP6 models, including moderate emissions constraints (rcp45), leading to a moderately hotter and drier future on average across CONUS, and no emissions constraints (rcp85), leading to a substantially hotter and drier future. These pathways are further combined with two CMIP model groups: one relatively cooler and one relatively hotter, based on their temperature sensitivity over CONUS. For each TGW scenario we compute daily potential evapotranspiration (PET), precipitation (P), mean temperature, heating degree days (HDD), cooling degree days (CDD), and a growing-season index (GSI). These daily fields are aggregated to monthly totals or means as appropriate and reprojected from the native TGW WRF grid to the Tethys 1/8$^{\circ}$ CONUS grid.

% COMMENT this should be moved to "temporal downscaling" subsection
The monthly water deficit used in irrigation temporal downscaling is defined as $\Delta_m = \max\!\big(\text{PET}_m - P_m,\, 0\big)$. The monthly GSI is derived from a simplified version of Jolly et al.\cite{Jolly_2005}, retaining the daily minimum-temperature and daylength indicators but omitting the vapour-pressure-deficit term. Specifically, for each day $d$,
\begin{equation}
f(T_{\min,d}) = \min\!\left(\max\!\left(\frac{T_{\min,d}+2}{7},\,0\right),\,1\right), \qquad
g(L_d) = \min\!\left(\max\!\left(L_d - 10,\,0\right),\,1\right),
\label{eq:gsi-components}
\end{equation}
where $T_{\min,d}$ is the daily minimum air temperature (\textdegree C) and $L_d$ is the daylength (hours). Monthly GSI is the daily-mean product over month $m$,
\begin{equation}
G_m = \big\langle f(T_{\min,d}) \cdot g(L_d) \big\rangle_{d \in m}.
\label{eq:gsi-monthly}
\end{equation}
Preprocessing code is archived under \texttt{scripts/0\_preprocessing/gsi\_nersc/} in the integration meta-repository (see ``Code availability'').

\subsection*{Spatial downscaling}

% COMMENT Why the pi? Probably a good idea to replace 
The location of demand within a region is assumed to follow that of a gridded proxy appropriate to the sector. For any sector with total demand $D_{\mathrm{region}}$ and proxy field $\pi_{\mathrm{cell}}$,
\begin{equation}
D_{\mathrm{cell}} = D_{\mathrm{region}} \times \frac{\pi_{\mathrm{cell}}}{\sum_{c \in \mathrm{region}} \pi_{c}}.
\label{eq:spatial}
\end{equation}
The following subsections specify the proxy used for each sector.

\subsubsection*{Irrigation}

GCAM-USA reports water demand for 13 crop classes: Corn, Wheat, Rice, RootTuber, OilCrop, SugarCrop, OtherGrain, FiberCrop, FodderGrass, FodderHerb, biomass, MiscCrop, and PalmFruit. We use Demeter\cite{Vernon-2018}, a land-use spatial disaggregation model, to produce scenario-consistent annual per-crop irrigated-area maps at 1/8$^{\circ}$. Demeter applies transition-priority rules that distinguish intensification (increase of a land type within a cell) from extensification (spread from an adjacent cell) to reconcile baseline high-resolution land use with GCAM-USA boundary conditions. Per-crop irrigated-area fields then serve as the proxy $\pi_{\mathrm{cell}}$ for each crop's irrigation demand via Eq.~(\ref{eq:spatial}).

\subsubsection*{Electricity}

Electricity-sector demand in this dataset represents cooling water consumed by thermoelectric generation (coal, natural gas, oil, nuclear, biomass, geothermal, concentrated solar thermal), excluding hydropower. Previous gridded products often used population as a proxy, on the reasoning that thermoelectric plants are sited near load centers. 
% COMMENT This should be introduced in the first section to  justify the proposed enhancements in Tethys 2.0
This assumption breaks down at 1/8$^{\circ}$: plants are sparser than people and are sited by a combination of cooling-water access, fuel logistics, transmission proximity, and siting constraints, not by local population. We therefore use the Capacity Expansion Regional Feasibility (CERF)\cite{Vernon2021} model to site projected power plants at $\approx$1-km resolution and 5-year increments, consistent with the installed-capacity trajectories output by GCAM-USA at state level. For each technology class, CERF-sited plants are aggregated to 1/8$^{\circ}$ grid cells and the cell-level installed capacity (MW) is used as the proxy $\pi_{\mathrm{cell}}$. State-level thermoelectric demand from GCAM-USA is then allocated proportionally to installed capacity within the state. For the historical period we use the 2015 plant inventory from the Global Power Plant Database v1.3 (GPPD) augmented with the IM3 experiment B CONUS plant inventory, merged onto the same 1/8$^{\circ}$ grid.

\subsubsection*{Municipal (public supply and domestic)}

SSP-consistent gridded population maps from Jones and O'Neill\cite{Jones_2016} are available at 10-year intervals between 2010 and 2100. For years in which population is required but unavailable, we linearly interpolate between the two bracketing decadal maps. Within each state, municipal demand from GCAM-USA is distributed to cells proportionally to projected population via Eq.~(\ref{eq:spatial}).

\subsubsection*{Livestock}

GCAM represents five livestock-product sectors (Beef, Dairy, Pork, Poultry, SheepGoat). Gridded livestock head-count maps from the GLW~3 dataset\cite{Gilbert2018} provide the spatial proxy. Mapping of GCAM sectors to GLW animals follows Table~\ref{tab:livestock}. GLW is published at 1/12$^{\circ}$ resolution for reference year 2010; we re-grid to 1/8$^{\circ}$ by allocating source-cell counts to target cells in proportion to fractional overlap, and hold the 2010 distribution fixed across all years and scenarios. The overall water-demand contribution from livestock is small (\textasciitilde2\% of CONUS total) so the impact of this static-distribution assumption on the aggregate record is limited; see ``Limitations''.

\begin{table}[ht]
\centering
\begin{tabular}{ll}
\toprule
\textbf{GCAM livestock sector} & \textbf{GLW proxy animal(s)} \\
\midrule
Beef      & Buffalo + Cattle \\
Dairy     & Buffalo + Cattle \\
Pork      & Pig \\
Poultry   & Chicken + Duck \\
SheepGoat & Sheep + Goat \\
\bottomrule
\end{tabular}
\caption{Mapping from GCAM livestock sectors to animal species in the GLW~3 dataset.}
\label{tab:livestock}
\end{table}

\subsubsection*{Manufacturing and mining}

For manufacturing and mining, population is used as the spatial proxy, matching the approach of prior Tethys work\cite{Khan2023}. Total U.S.~demand for these two sectors is small compared to irrigation or thermoelectric, which limits the effect of this less-refined proxy on the aggregate record. See ``Limitations'' for localized implications.

\subsection*{Temporal downscaling}

GCAM-USA operates at 5-year steps. Spatially downscaled outputs are linearly interpolated in time to annual resolution before monthly downscaling. Sector-specific formulas then distribute the annual demand at each cell into 12 monthly values. Livestock, manufacturing, and mining monthly demand is assumed uniform at 1/12 of the annual value; irrigation, electricity, and domestic follow climate-driven formulations described below.

\subsubsection*{Irrigation}

For each cell and year, we compute a monthly weight field from the growing-season index $G_m$ (Eq.~\ref{eq:gsi-monthly}), the monthly deficit $\Delta_m$, and month length $N_m$:
\begin{equation}
\tilde{w}_m = \frac{\Delta_m \, G_m}{N_m}\quad, \qquad w_m = \frac{\tilde{w}_m}{\sum_{k=1}^{12}\tilde{w}_k}
\label{eq:irr-weights}
\end{equation}
so that $\sum_m w_m = 1$. Monthly irrigation demand is then $D_m = w_m \, D_{\mathrm{year}}$. Equation~(\ref{eq:irr-weights}) operationalizes the combined deficit-and-growing-season approach of Moore et al.\cite{Moore_2015} and Roy et al.\cite{Roy_2005}: water demand is concentrated in months that are simultaneously water-limited ($\Delta_m > 0$) and vegetatively active ($G_m > 0$). The weights are per-cell and differ across cells and years.

\subsubsection*{Electricity}

Monthly electricity water demand is distributed according to the share of annual electricity use that services heating, cooling, and all other uses, each modulated by a temperature-dependent distribution. Region-level shares of annual electricity consumption for heating $p_{\mathrm{heat}}$, cooling $p_{\mathrm{cool}}$, and other $p_{\mathrm{other}}$ come from GCAM-USA. At each cell, define annual sums $H_y = \sum_m \text{HDD}_m$ and $C_y = \sum_m \text{CDD}_m$ and monthly distribution fields $\hat{h}_m$ and $\hat{c}_m$ with
\begin{equation}
(\hat{h}_m,\,\hat{c}_m) =
\begin{cases}
\big(\text{HDD}_m / H_y,\ \text{CDD}_m / C_y\big)      & H_y \ge 650 \text{ and } C_y \ge 450 \\[2pt]
\big(\text{HDD}_m / H_y,\ \text{HDD}_m / H_y\big)      & H_y \ge 650 \text{ and } C_y < 450 \\[2pt]
\big(\text{CDD}_m / C_y,\ \text{CDD}_m / C_y\big)      & H_y < 650 \text{ and } C_y \ge 450 \\[2pt]
(1/12,\ 1/12)                                          & H_y < 650 \text{ and } C_y < 450
\end{cases}
\label{eq:elec-cases}
\end{equation}
with $\hat{o}_m = 1/12$ always. Monthly demand is
\begin{equation}
D_m = D_{\mathrm{year}} \times \big(p_{\mathrm{heat}}\,\hat{h}_m + p_{\mathrm{cool}}\,\hat{c}_m + p_{\mathrm{other}}\,\hat{o}_m\big).
\label{eq:elec-demand}
\end{equation}
The first case in Eq.~(\ref{eq:elec-cases}) applies to cells with both a meaningful heating season (annual HDD $\ge 650$) and a meaningful cooling season (annual CDD $\ge 450$). The second and third cases collapse onto whichever signal is non-trivial, which prevents division by near-zero annual sums in climates dominated by one extreme. The fourth case reverts to a uniform distribution where neither threshold is exceeded. The threshold convention follows Huang et al.\cite{hess-22-2117-2018}.

\subsubsection*{Domestic}

For the domestic (residential public-supply) sector we use the temperature-anomaly formula of Wada et al.\cite{Wada2011}:
\begin{equation}
D_m = \frac{D_{\mathrm{year}}}{12} \times \left(\frac{T_m - \bar{T}}{T_{\max} - T_{\min}}\,R + 1\right),
\label{eq:dom-monthly}
\end{equation}
where $T_m$ is the monthly mean temperature at the cell, $\bar{T}$ is the annual mean temperature, $T_{\max}$ and $T_{\min}$ are the annual extremes, and $R$ is a region-scale amplitude coefficient. The amplitude $R$ is taken from a static regional map and reflects the sensitivity of observed domestic use to seasonal temperature swings; higher $R$ implies stronger seasonal amplification.

\subsection*{Gridded renewable vs.\ non-renewable source shares}

GCAM-USA solves for the share of each basin's withdrawals met from renewable (surface) versus non-renewable (groundwater) supply by cost-based allocation. Mapping that basin-scale split onto the 1/8$^{\circ}$ grid requires two steps. First, the basin-level share is applied uniformly to all end-use sectors in the basin, except that thermoelectric demand is assumed to use surface water only and the remaining sectoral shares are renormalized accordingly. Second, we anchor the historical spatial pattern of renewable/non-renewable mix to observed USGS data by rescaling each cell's GCAM share against its historical-2015 value while preserving the USGS baseline spatial pattern. Specifically, let $s^{\mathrm{GCAM}}_{c,y}$ denote the GCAM-derived renewable share at cell $c$ in year $y$, and $s^{\mathrm{USGS}}_c$ the static 2015 USGS-derived share at cell $c$. For the subset of cells $\mathcal{M}$ for which both $s^{\mathrm{GCAM}}_{c,\,2015}>0$ and $s^{\mathrm{USGS}}_c$ is available, the adjusted share is
\begin{equation}
s^{\mathrm{adj}}_{c,y} =
\begin{cases}
\displaystyle \min\!\left(1,\; s^{\mathrm{USGS}}_c \times \frac{s^{\mathrm{GCAM}}_{c,y}}{s^{\mathrm{GCAM}}_{c,\,2015}}\right) & c \in \mathcal{M} \\[10pt]
s^{\mathrm{GCAM}}_{c,y}                                                                                        & c \notin \mathcal{M}
\end{cases}
\label{eq:source-shares}
\end{equation}
The $\min(\cdot,1)$ clip bounds ratio amplifications where the GCAM 2015 baseline is small. The USGS spatial pattern is thus anchored in place where observable, and the GCAM scenario trajectory enters only as a per-cell temporal ratio.

\subsection*{Future projection methods}

The dataset combines simulation outputs that were produced, in some cases, by independent teams using scenario assumptions that are consistent but not strictly identical. In particular: the GCAM-USA run uses its own scenario configuration for each SSP$\times$RCP combination; Demeter is run with the corresponding GCAM scenario to produce per-crop irrigated-area maps; CERF sites power plants using GCAM-USA state-level capacity trajectories; the Jones and O'Neill SSP3/SSP5 population maps are used as-is; and the TGW-WRF climate sample is drawn from a fixed ensemble conditioned on the RCP. We do not attempt to perfectly co-calibrate these inputs, which would be prohibitively expensive; instead, we document the specific versions used (see ``Code availability'') and treat the resulting record as scenario-plausible but not a single self-consistent Earth-system simulation. This ``frankenstein'' design is an established compromise in multi-sector dynamics modeling\cite{Khan2023} and is, in our view, the right one for a dataset intended to support sensitivity and adaptation studies rather than causal attribution.

% -----------------------------------------------------------------------------
\section*{Data Records}

% COMMENT replace MSD-Live link with DOI before publication
The dataset is openly available for public access at MSD-Live (\url{https://data.msdlive.org/uploads/p4xce-e8822}) and the Tethys model is available at: \href{https://github.com/JGCRI/tethys}{github.com/JGCRI/tethys}.

The dataset is openly available and  permanently on MSD-Live (\url{https://data.msdlive.org/uploads/p4xce-e8822}). Scenario directories (Table~\ref{tab:scenarios}) each contain per-sector netCDF~4 files following the naming convention
\texttt{<Sector>\_<demand\_type>[\_monthly].nc}, where \texttt{<Sector>} is one of \texttt{Domestic}, \texttt{Electricity}, \texttt{Irrigation}, \texttt{Livestock}, \texttt{Manufacturing}, \texttt{Mining}; \texttt{<demand\_type>} is \texttt{withdrawals} or \texttt{consumption}; and the \texttt{\_monthly} suffix distinguishes monthly files from their annual counterparts. For irrigation, an additional \texttt{\_with\_losses} suffix marks files that include conveyance losses.

Each scenario directory also contains \texttt{gridded\_runoff\_shares.nc} (per-year, per-cell renewable share $s^{\mathrm{adj}}_{c,y}$ from Eq.~\ref{eq:source-shares}) and two YAML configuration files (\texttt{config\_withdrawals.yaml}, \texttt{config\_consumption.yaml}) that record the exact Tethys run configuration used to produce the files, for reproducibility. Figure~\ref{fig:data-listing} shows an example CDL listing. All sector files share the same \texttt{(year, lat, lon)} or \texttt{(year, lat, lon, month)} schema with sector-specific sub-variables.

% Need to revise this based on new runs
\begin{table}[ht]
\centering
\begin{tabular}{lll}
\toprule
\textbf{Scenario directory} & \textbf{Time span} & \textbf{Description} \\
\midrule
\texttt{historical}         & 1980--2019 & Observed-period baseline using historical climate and population. \\
\texttt{rcp45cooler\_ssp3}  & 2020--2100 & Lower-emissions, cooler TGW sample, SSP3 socioeconomics. \\
\texttt{rcp45cooler\_ssp5}  & 2020--2100 & Lower-emissions, cooler TGW sample, SSP5 socioeconomics. \\
\texttt{rcp45hotter\_ssp3}  & 2020--2100 & Lower-emissions, hotter TGW sample, SSP3. \\
\texttt{rcp45hotter\_ssp5}  & 2020--2100 & Lower-emissions, hotter TGW sample, SSP5. \\
\texttt{rcp85cooler\_ssp3}  & 2020--2100 & Higher-emissions, cooler TGW sample, SSP3. \\
\texttt{rcp85cooler\_ssp5}  & 2020--2100 & Higher-emissions, cooler TGW sample, SSP5. \\
\texttt{rcp85hotter\_ssp3}  & 2020--2100 & Higher-emissions, hotter TGW sample, SSP3. \\
\texttt{rcp85hotter\_ssp5}  & 2020--2100 & Higher-emissions, hotter TGW sample, SSP5. \\
\bottomrule
\end{tabular}
\caption{Scenario directories in the published dataset. Each directory contains per-sector withdrawal and consumption netCDFs, a gridded runoff-share file, and the two Tethys YAML configurations used to generate them.}
\label{tab:scenarios}
\end{table}

% COMMENT Too much info for the main text 
\begin{figure}[ht]
\centering
\begin{verbatim}
netcdf Livestock_withdrawals_monthly {
dimensions:
    year = 46 ;
    lat = 224 ;
    lon = 464 ;
    month = 12 ;
variables:
    float Beef(year, lat, lon, month) ;
    float Dairy(year, lat, lon, month) ;
    float Pork(year, lat, lon, month) ;
    float Poultry(year, lat, lon, month) ;
    float SheepGoat(year, lat, lon, month) ;
    double lat(lat) ;
    double lon(lon) ;
    int64 year(year) ;
    int64 month(month) ;
    int64 spatial_ref ;
}
\end{verbatim}
\caption{CDL excerpt of a monthly livestock-withdrawals netCDF. Other sectors share the same \texttt{(year, lat, lon [, month])} layout with sector-specific sub-variables.}
\label{fig:data-listing}
\end{figure}

% -----------------------------------------------------------------------------
\section*{Technical Validation}

We validate the downscaled dataset at the three sectors that together account for over 90\% of CONUS water demand: irrigation, thermoelectric, and domestic (public supply)\cite{skinnerWaterWithdrawalConsumption2025}. Reference data are the USGS five-year water-use records, refreshed in January~2025, at the HUC12 scale for public supply and irrigation and per-plant for thermoelectric. We proceed from coarse to fine and from total to seasonal: first comparing CONUS annual totals, then moving to HUC6 spatial agreement, then to the HUC6 annual scatter and correlation, and finally the CONUS monthly cycle. The goal is not to declare a ``true'' dataset, since neither USGS nor Tethys are direct observations at 1/8$^{\circ}$ resolution, but to establish that the downscaled record reproduces the dominant features of observed spatial and temporal demand patterns, with quantified bias where it departs.

\subsection*{CONUS annual totals}

Figure~\ref{fig:annual-total-timeseries} compares CONUS-aggregated annual demand for each of the three sectors and their sum. Tethys reproduces both the magnitude and long-term trend of USGS totals, within 10\% at annual resolution. The GCAM-USA 5-year time step manifests in the Tethys annual irrigation demands as reduced interannual variability relative to USGS. Magnitudes differ sharply across sectors, as expected: irrigation dominates consumptive use, while thermoelectric and irrigation are of comparable magnitude in withdrawals. The decline Electricity withdrawals reflects a known trend due in large part to the switch from coal-fired plants to other technologies \cite{skinnerWaterWithdrawalConsumption2025}. Figure~\ref{fig:annual-total-boxplot} shows the distribution of percent difference between the two records across the HUC6 regions. Domestic and Irrigation show sizable sector-level biases in opposing directions at the HUC6 scale, which partially cancel at the CONUS total.

\begin{figure}[ht]
\centering
\includegraphics[width=\linewidth]{val1-huc06-usgs-tethys-annual-total.png}
\caption{CONUS annual water use by sector from Tethys (black) and USGS (orange), 2000--2020.}
\label{fig:annual-total-timeseries}
\end{figure}

% COMMENT Could we jazz this up a bit. Also need to explain why we claimed 10% error in the abstract but these y axis spans -75 to 25%. also need to explain outliers 
\begin{figure}[ht]
\centering
\includegraphics[width=\linewidth]{val2-huc06-usgs-tethys-annual-total-pdiff.png}
\caption{Percent difference between Tethys and USGS annual water use across HUC6 regions, by sector and demand type.}
\label{fig:annual-total-boxplot}
\end{figure}

\subsection*{HUC6 spatial agreement}

Figure~\ref{fig:map-pbias} maps the percent difference in annual average use at HUC6 scale. Tethys tends to exceed USGS in the Electricity sector, especially in the eastern and southeastern U.S.~where CERF-sited plants concentrate cooling demand in basins where USGS reports lower per-HUC totals. Tethys falls below USGS in Irrigation in several western basins --- notably the lower Colorado and parts of the Central Valley --- where the GSI-weighted monthly distribution underweights months that USGS records as heavy irrigation under observed 2015 conditions. At basin-wise correlation (Figure~\ref{fig:huc-correlation}), annual-average agreement at HUC6 ranges from 0.71 (Domestic) to 0.95 (Irrigation), which we take as evidence that the spatial pattern of use is well captured even where the magnitude is biased.

\begin{figure}[ht]
\centering
\includegraphics[width=\linewidth]{val3-huc06-pdiff-usgs-tethys.png}
\caption{Percent difference in annual-average HUC6 water use, Tethys minus USGS, for Domestic, Electricity, and Irrigation sectors.}
\label{fig:map-pbias}
\end{figure}

% COMMENT Could we please facet_grid here and maybe some styling or coloring of scatter by something meaningful 
\begin{figure}[ht]
\centering
\includegraphics[width=\linewidth]{val4-huc06-scatter-usgs-tethys.png}
% COMMENT Need to mention the year.
\caption{HUC6 annual average demand, Tethys versus USGS, by sector and demand type. Pearson correlations range from 0.71 (Domestic) to 0.95 (Irrigation).}
\label{fig:huc-correlation}
\end{figure}

\subsection*{Seasonal cycle}

Figure~\ref{fig:monthly} compares the mean monthly cycle at CONUS scale. Irrigation and Electricity withdrawals reproduce the seasonal shape observed by USGS closely. Electricity consumption shows a consistent offset in non-summer months, likely from the GCAM-USA representation of non-cooling electricity water use. Domestic shows a broadly consistent positive bias across months; this reflects the $R$ amplitude coefficient used in Eq.~(\ref{eq:dom-monthly}), which was calibrated to aggregate USGS demand rather than to the post-2015 public-supply subset.

\begin{figure}[ht]
\centering
\includegraphics[width=\linewidth]{val5-huc06-usgs-tethys-monthly-total.png}
\caption{Mean monthly CONUS water use by sector, Tethys (black) and USGS (orange), averaged over 2000--2020.}
\label{fig:monthly}
\end{figure}

\subsection*{Inter-scenario consistency}

Figure~\ref{fig:scenarios} shows CONUS annual demand across all nine scenarios for the three largest sectors and their sum. The scenario spread exhibits the expected signature of each input: strong SSP-driven divergence in Domestic (SSP5 rising with population, SSP3 declining), large inter-scenario variability in Irrigation driven by RCP, and a monotone decline in Electricity across all futures that reflects the GCAM-USA projection of declining thermoelectric water use under continued fleet turnover. The discontinuity at 2020 between the historical and future lines reflects the switch from the ERA5‑based reanalysis historical run to the TGW-driven future simulations that restart the historical weather sequence with added thermodynamic warming; within the future period the scenarios evolve smoothly and consistently with their  drivers.

\begin{figure}[ht]
\centering
\includegraphics[width=\linewidth]{val6-scenarios-annual-conus-timeseries.png}
\caption{Annual CONUS water withdrawals by sector and scenario. Historical simulation (black, 1975--2015) is shown alongside eight future scenarios (2020--2100) that vary by RCP (color), climate-model sample (line width), and SSP (line style). Each panel uses an independent $y$-axis.}
\label{fig:scenarios}
\end{figure}

\subsection*{Limitations}

Several known limitations should inform reuse of the dataset:
\begin{itemize}
    \item \textbf{Livestock spatial stationarity.} The GLW~3 livestock distribution is held fixed at 2010 across all years and scenarios. Livestock is \textasciitilde2\% of CONUS demand, but within individual basins --- particularly in states with rapidly shifting livestock composition --- this may introduce local bias.
    \item \textbf{Manufacturing and mining proxy.} Using population as the proxy for these two sectors captures average patterns of labor-associated demand but misses site-specific industrial or mining facilities whose water withdrawals can dominate local basin totals.
    \item \textbf{Electricity source-share carve-out.} Eq.~(\ref{eq:source-shares}) assigns thermoelectric demand to surface water only and re-normalizes remaining-sector shares. Where a thermoelectric plant is served from a local groundwater source (rare but documented), this assumption places the withdrawal on surface water and slightly overstates surface demand elsewhere.
    \item \textbf{Riparian vs.\ reservoir allocation.} The dataset represents demand at the cell where use occurs, which is not necessarily the cell from which water is withdrawn in reservoir-fed or interbasin-transfer systems. Downstream routing and water-management models (e.g., mosartwmpy) handle the supply-side re-allocation.
    \item \textbf{Simplified GSI.} The growing-season index used for irrigation temporal downscaling (Eqs.~\ref{eq:gsi-components}--\ref{eq:gsi-monthly}) drops the vapour-pressure-deficit (VPD) component of the original Jolly et al.\cite{Jolly_2005} formulation. This simplification is consistent with Moore et al.\cite{Moore_2015} but may underweight drought months in humid climates where VPD is the binding constraint.
    \item \textbf{Conveyance losses.} Irrigation files exist with and without conveyance losses applied (\texttt{\_with\_losses} suffix). Users coupling the dataset to hydrologic routing should select the appropriate variant for their application.
\end{itemize}

% -----------------------------------------------------------------------------
\section*{Usage Notes}

% COMMENT Too much information for a paper, this can be moved to the README for the dataset
The dataset is provided in netCDF~4 and is readable in any standard scientific-computing environment. A typical workflow in Python:
\begin{verbatim}
import xarray as xr
ds = xr.open_dataset(
    "rcp45cooler_ssp3/Domestic_withdrawals_monthly.nc"
)
# CONUS total by year and month:
total = ds["Domestic"].sum(dim=("lat", "lon"))
\end{verbatim}
Conversion between the native unit (km$^3$~yr$^{-1}$) and the USGS unit (million gallons per day, MGD): $1\text{ km}^3\text{ yr}^{-1} \approx 723.8\text{ MGD}$ (using $\text{km}^3$\,yr$^{-1}\rightarrow\text{MGD}$ $= 264\,172.05124/365$).

HUC-level aggregation is supported in the companion integration meta-repository under \texttt{validation/1-postprocess-tethys.py}, which uses the \texttt{xagg} library to compute overlap-weighted aggregates from the 1/8$^{\circ}$ grid to HUC2, HUC4, HUC6, HUC8, or HUC12 polygons.

% COMMENT We should lead with this figure as Fig 1
\begin{figure}[ht]
\centering
\includegraphics[width=\linewidth]{usage1-dominant-sector-tethys-grid.png}
\caption{Dominant water-use sector at each 1/8$^{\circ}$ cell, by annual-average consumption.}
\label{fig:dominant-sector}
\end{figure}


% COMMENT this section repeats some info already in the Methods. review it again and consider where trimming may make sense.
\section*{Improvements over previous version}

Compared with the prior Tethys global product\cite{Khan2023}, the dataset presented here advances the representation of demand in six specific ways.

\paragraph{GCAM-USA integration.} The prior record used the global GCAM configuration, which disaggregated U.S.~demand from a single national total. This dataset uses GCAM-USA, which resolves U.S.~electricity generation, manufacturing, and municipal demand at the state level and crop irrigation demand at state$\,\times\,$basin intersections. State-level inputs eliminate the artifact of national-average demand being pushed uniformly into states with very different sectoral mixes, and they allow the subsequent spatial downscaling to respect state-level regulatory and policy boundaries.

\paragraph{CERF-based power-plant siting.} Prior gridded electricity demand used population as the spatial proxy, which approximates the location of load rather than the location of cooling demand. This dataset uses explicit power-plant siting from the CERF model\cite{Vernon2021}, which places thermoelectric plants at $\approx$1-km resolution based on siting feasibility (exclusion layers, transmission proximity, cooling-water access) and then aggregates installed capacity to the 1/8$^{\circ}$ grid as the proxy. The result is a markedly more realistic spatial distribution of thermoelectric demand, particularly in regions where load centers and generation sites are geographically decoupled.

\paragraph{SSP-consistent population.} Prior work used a static base-year population map to distribute municipal, manufacturing, and mining demand across all years, including future projections. This dataset uses SSP3 and SSP5 gridded population projections from Jones and O'Neill\cite{Jones_2016}, linearly interpolated from their decadal native resolution to annual values. Scenario-consistent population is essential for the SSP-differentiated futures: SSP3 and SSP5 diverge markedly in U.S.~population growth, and that divergence propagates into the public-supply demand field in a way that a static map cannot represent.

\paragraph{GSI-based irrigation temporal downscaling.} Prior irrigation monthly weights used a simpler temperature-driven scheme. This dataset computes per-cell monthly weights (Eq.~\ref{eq:irr-weights}) from the TGW-WRF monthly deficit and growing-season index, ensuring that the monthly distribution of irrigation demand responds to climate-consistent drought and growing-condition signals rather than to a static seasonal template. In consequence, the monthly irrigation distribution varies from year to year and across scenarios, as it should.

\paragraph{USGS-anchored source-share attribution.} Prior work applied GCAM's basin-level renewable vs.\ non-renewable split uniformly across cells within each basin. This dataset uses Eq.~(\ref{eq:source-shares}) to anchor the historical 2015 spatial pattern to USGS observations while preserving the temporal evolution from GCAM. In regions with observable baseline data this reduces bias materially; outside those regions, the GCAM share is preserved transparently.

\paragraph{Resolution refinement.} Spatial resolution is refined from 1/2$^{\circ}$ to 1/8$^{\circ}$, a factor-of-four improvement in each dimension and 16x improvement in areal resolution. Combined with the finer proxy data (CERF plants at \textasciitilde1~km, Demeter LULC at 1/8$^{\circ}$, Jones population at 1/8$^{\circ}$), the 1/8$^{\circ}$ grid allows the resulting demand field to represent sub-state variation relevant to river routing, reservoir management, and local planning applications.

Together, these six advances produce a dataset that validates well against USGS at HUC6 (annual correlations of 0.71--0.95) and that supports detailed scenario. In addition this dataset is validated against the refreshed January 2025 USGS water-use record, which includes updated thermoelectric cooling-water estimates\cite{skinnerWaterWithdrawalConsumption2025}.

% -----------------------------------------------------------------------------
\section*{Code availability}

All code used to generate and validate the dataset is publicly available under permissive open-source licenses.
\begin{itemize}
    \item \textbf{Tethys downscaling package}: \url{https://github.com/JGCRI/tethys} (\texttt{pip install tethys-downscale}). The specific version used for the published dataset is pinned in the scenario YAML config files included with each scenario directory.
    \item \textbf{Integration meta-repository}: \url{https://github.com/IMMM-SFA/tethys_integration_metarepo}. Contains the per-scenario run drivers (\texttt{scripts/1\_runs/im3\_tethys\_runs/}), climate-forcing preprocessing (\texttt{scripts/0\_preprocessing/gsi\_nersc/} and \texttt{scripts/0\_preprocessing/compute\_\{gsi,deficit,monthly\_weights\}.py}), runoff-share adjustment (\texttt{scripts/2\_postprocess/adjust\_runoff\_shares/}), and the numbered validation pipeline (\texttt{validation/}).
    \item \textbf{Demeter} (land-use downscaling, used to produce per-crop irrigated-area proxies): \url{https://github.com/IMMM-SFA/demeter}\cite{Vernon-2018}.
    \item \textbf{CERF} (power-plant siting): \url{https://github.com/IMMM-SFA/cerf}\cite{Vernon2021}.
    \item \textbf{TGW-WRF climate forcing}\cite{Jones2023TGW}: \url{https://tgw-data.msdlive.org/}.
\end{itemize}
% COMMENT this referes to a specific file in the metarepo, please only refer to the repo in general to abv
To reproduce the published record for one scenario, follow the \texttt{PIPELINE.md} in the integration meta-repository, which documents required inputs, expected outputs, and the command sequence for each stage.

% -----------------------------------------------------------------------------
\section*{Acknowledgments}

This research was supported by the U.S.~Department of Energy, Office of Science, as part of research in MultiSector Dynamics, Earth and Environmental System Modeling Program. We thank the TGW-WRF team for the downscaled climate product, the GCAM-USA team for the scenario simulations, and the USGS water-use program for the reference observations used in validation.

\section*{Author contributions statement}

% TODO: coauthors to fill in per CRediT roles. Illustrative template:
C.B.~Coordinated the downscaling runs, conducted the HUC-scale validation, and drafted this version of the manuscript based on an earlier draft by I.T.
H.N.~Developed and maintained the Tethys model, conducted the runs, and helped with the interpretation of GCAM-USA and Tethys outputs.
T.T.~Extended the Tethys package to support GCAM-USA inputs, developed the first version of the Tethys metarepo, and carried out the initial scenario runs.
I.T.~Wrote initial draft of this manuscript, contributed to the CERF--Tethys integration and power-plant proxy construction.
H.E.~Produced the TGW-WRF preprocessing pipeline for PET, HDD/CDD, and GSI.
K.M.~Developed CERF CERF--Tethys integration and power-plant proxy construction.
N.V.~Contributed to scenario design and interpretation.
N.S.~Project management, scenario design and interpretation.
J.R.~Provided overall project guidance and scientific review.

\section*{Competing interests}

The authors declare no competing interests.

\bibliography{Tethys}

\end{document}
